\chapter{Commercial System Survey}
\label{ch:survey}
% Retain the historical source catalog without using these records as evidence
% for removed universal performance, algorithm or feature claims.
\nocite{fccTman4,flightscopebgh,foresighthmt,garminr10accuracy,garminr10data,mevoteardown,mygolfspymlm2pro,uneekordimple,uneekorphotometric,us12401909,playbetterblp}

A product survey is useful only if its categories preserve what the instrument observes, what its estimator assumes and what the purchased configuration reports.
\Cref{tab:survey} summarizes selected architectures using manufacturer documentation checked on October 4, 2026.
It is neither an accuracy ranking nor an exhaustive list of current products.
Firmware, enabled software, ball type, club markers and installation geometry can change the available outputs.

The measurement hierarchy in \cref{tab:hierarchy} applies to every row.
Camera images and radar signals are observations; launch and club states are estimates from those observations; ratios and coordinate transformations are derived quantities; simulated downrange flight is a model prediction.
An output advertised as ``measured'' still requires a definition, an estimator and uncertainty assessment.
Conversely, a radar-based face-angle output cannot be assigned to a particular D-plane inverse merely from its sensor category.
Where public documentation does not disclose that algorithm, the table does not invent one.

\begingroup
\small
\setlength{\tabcolsep}{4pt}
\begin{longtable}{@{}p{3.0cm}p{5.2cm}p{6.2cm}@{}}
\caption{Documented Launch-Monitor Architectures and Interpretation Limits}\label{tab:survey}\\
\toprule
\textbf{System} & \textbf{Observation Architecture} & \textbf{Configuration and Interpretation} \\
\midrule
\endfirsthead
\toprule
\textbf{System} & \textbf{Observation Architecture} & \textbf{Configuration and Interpretation} \\
\midrule
\endhead
\bottomrule
\endfoot
TrackMan 4 & Synchronized short- and long-range radar with a camera & Optical/radar integration; indoor distance is specified from unit to net, not wholly as ball flight. \\
TrackMan iO & Radar, high-speed camera and alignment camera & Overhead indoor geometry; club outputs depend on the software package. \\
FlightScope X3 / X3C & Radar and camera Fusion Tracking & Model-specific setup; X3 documentation includes face impact, and current X3C documentation also lists it. \\
FlightScope Mevo+ & Radar with a built-in camera & Pro and Face Impact Location are software options; neither means adding a previously absent camera. \\
Garmin R10 & Radar & Standard-ball spin requirements differ from RCT-assisted indoor requirements; calculated face angle is not a direct image observation. \\
Garmin R50 & Three cameras & Club sticker required; an impact replay is not by itself evidence of calibrated numerical impact coordinates. \\
Full Swing KIT & Radar with a video camera & Manufacturer describes ML-enhanced radar; the public account does not establish a camera contribution to every estimate. \\
Rapsodo MLM2PRO & Radar and two cameras & RPT optical-pattern balls required for reported spin; RCT is a different technology. \\
Foresight GCQuad / QuadMAX / GCHawk & Four-camera family & Club-marker mode and enabled features matter; do not transfer every model's feature list to the entire family. \\
Foresight GC3 / Bushnell Launch Pro & Three-camera family & GC3's documented club set is speed, path, attack and smash; face orientation and impact location are not established by that list. \\
Uneekor EYE XO2 & Three overhead cameras & The manual requires club stickers for club data; overhead placement does not imply marker-free operation. \\
Uneekor EYE MINI & Two cameras & Club Optix supplies impact video; manual/support and product-page feature descriptions need model/version reconciliation. \\
ProTee VX & Two overhead cameras & Vendor lists sticker-free ball and club estimates, including numerical impact coordinates; these are capability claims. \\
SkyTrak ST+ & Photometric ball tracking with dual-radar club tracking & Manufacturer describes learned club estimation; do not assume every ST MAX specification is identical. \\
OpenFlight & Radar development architecture & Hardware interfaces and experimental processing are discussed in \cref{app:impl}; no universal spin-detection rate follows. \\
PiTrac & Camera and infrared-strobe development architecture & Configuration and calibration are version-dependent; low-cost parts do not establish validated accuracy. \\
\end{longtable}
\endgroup

\section{Radar and Hybrid Systems}

\subsection{TrackMan 4 and iO}

TrackMan describes two synchronized radar subsystems and optical assistance for TrackMan~4~\cite{trackmanspecs,trackmanoert}.
Its published minimum indoor distance of 4.7\,m is measured from the device to the net; part of that distance lies behind the ball.
Treating the entire distance as a required ball-flight segment changes the installation claim.
The overhead iO has a different geometry and lists 24\,GHz radar, a high-speed camera capable of up to 4,600 frames/s and an alignment camera~\cite{trackmanspecs}.
It is not simply a rear-mounted TrackMan~4 repackaged above the golfer.
The listed Home package club output is narrower than the Home Complete/Commercial set, which includes face orientation and impact location.
Room-clearance and capture requirements still apply even when a product does not require a long radar flight segment.
These specifications establish advertised configuration, not common error bounds between packages or measurement equivalence with other brands.

\subsection{FlightScope X3, X3C and Mevo+}

FlightScope's Fusion Tracking combines radar and image information; Mevo+ already contains a camera used for alignment, video and tracking~\cite{flightscopefaq}.
The Pro Package enables additional software outputs, rather than converting a camera-free unit into a camera-equipped one.
The separate Face Impact Location option is available for Mevo+ without requiring Pro; FlightScope's X3 account describes impact location as included, and its X3C page also lists that feature~\cite{flightscopefaceimpact,flightscopex3csetup}.
Consequently, a categorical ``no impact location'' entry for these systems is incorrect.
Short-flight spin operation depends on the prescribed ball treatment, spacing and environment; face-impact imaging also has its own lighting requirements.
Published patents explain possible sensing methods, but they do not prove that each current model implements every claimed embodiment.
The distinction between patent evidence and deployed-product evidence is developed in \cref{ch:patents}.

\subsection{Garmin Approach R10}

Garmin lists club-head speed, ball speed, launch angle and launch direction as radar-measured quantities and identifies face angle as calculated~\cite{garminr10support}.
That statement does not identify all details of the internal estimator, nor does it make the remaining displayed metrics independent observations.
Its standard-ball guidance ties measured spin to substantial flight distance and ball speed; those conditions should not be imposed unchanged on its separate RCT-assisted indoor mode.
The RCT guidance specifies a minimum ball-to-screen distance and sufficient observed rotation, with longer flight sometimes required; an italic spin result indicates estimation~\cite{garminr10rct}.
Titleist's Radar Capture Technology (RCT) uses an embedded radar-reflective structure, unlike the optical RPT pattern.
A vendor tolerance or signal-detection percentage also cannot substitute for an independent, condition-specific error distribution.

\subsection{Full Swing KIT and Rapsodo MLM2PRO}

Full Swing describes KIT as an ML-enhanced radar system with 16 reported data points and video capability~\cite{fullswingkit}.
The presence of a camera does not establish that its images enter every kinematic estimate.
The cited manufacturer technology page does not substantiate a general third-party finding of 1--2\% agreement with other launch monitors, so no such accuracy claim is made here.

Rapsodo's MLM2PRO combines radar with two cameras and requires its RPT-pattern balls to report spin rate and spin axis~\cite{rapsodofaq,rapsodomlm2prorpt}.
RPT denotes Rapsodo Precision Technology; it is not the Titleist RCT radar treatment.
The current FAQ lists 15 metrics, including club path and attack angle, but does not list face angle.
That feature list therefore cannot support calling the device a demonstrated face-angle estimator.
Its impact video, launch estimates and simulated flight remain distinct outputs.
Neither a marketing comparison percentage nor an unsupported camera frame rate establishes accuracy across balls, clubs and capture conditions.

\section{Camera Systems}

\subsection{Foresight Families}

The GCQuad manual distinguishes a single-marker club mode from a four-marker mode that adds orientation information; its driver and iron diagrams show a noncollinear layout~\cite{foresightquadmanual}.
Supported GCQuad-family features include face and impact information, but QuadMAX and overhead GCHawk capabilities should be checked against their own model documentation and enabled features~\cite{foresightmarkers}.
No generic four-camera description establishes the exact spin algorithm, image count per shot or every club metric.
The robot comparisons in \cref{ch:accuracy} concern the tested conditions; they do not establish a universal class-leading accuracy ranking.

GC3 and Bushnell Launch Pro belong to the three-camera family, but branding and subscriptions should not replace model-specific entitlement checks~\cite{foresightgc3,bushnelllaunchpro}.
The GC3 product list specifies club speed, path, attack angle and smash factor.
Those outputs do not include the GCQuad's full face-orientation and impact-location set.
The point-velocity and marker-registration distinctions in \cref{sec:camclub} apply even when two units display the same metric name.

\subsection{Uneekor and ProTee}

The EYE~XO2 manual describes three cameras and club stickers; its mounting and calibration instructions are part of the measurement configuration~\cite{uneekorxo2manual}.
Its optical ball tracking does not make its club tracking sticker-free.
EYE~MINI is a different, portable two-camera model with impact video.
Uneekor's model-specific support requires club stickers, whereas its current product page calls them optional or recommended and displays a broader club-metric list~\cite{uneekorminisupport,uneekoreyexo}.
That documentation discrepancy should be resolved for the installed hardware/software before relying on numerical face or impact outputs; video replay alone is insufficient evidence.

ProTee VX's current product documentation explicitly advertises two cameras, sticker-free ball and club tracking, and numerical horizontal and vertical impact coordinates~\cite{proteevx}.
That is evidence of an advertised feature set, not an independent validation of the underlying image estimator.
Overhead placement, learned image interpretation and a large hitting zone each solve different parts of the observation problem.
None removes the need for calibration, adequate visibility and a declared face-coordinate convention.

\subsection{Garmin R50 and SkyTrak ST+}

Garmin documents three cameras for the R50 and requires a club-face sticker for club tracking~\cite{garminr50,garminr50support,garminr50stickers}.
The manual's impact-viewing functions should be distinguished from a verified specification of numerical impact coordinates.
Its use of cameras instead of the R10's radar is an architectural difference, not by itself proof of a quantified accuracy improvement.

SkyTrak describes photometric ball tracking and dual-radar club estimation with machine learning for ST+~\cite{skytrakplus}.
This is another reason to avoid a strict camera-versus-radar taxonomy.
The existence of another product, ST MAX, does not justify pooling its hardware, software and feature specifications with ST+ without separate documentation.

\section{Open-Source Systems and Comparison Practice}

The OpenFlight-oriented hardware example and PiTrac illustrate radar and strobed-camera development routes~\cite{an029,pitrac}.
The sensor interfaces, sample cadence and observability limits in \cref{app:impl} determine what can be inferred from available data; an experimental spin-detection rate is not a transferable product specification.
PiTrac's evolving documentation likewise needs a pinned build and calibration record before numerical performance can be reproduced.
A subtotal for computers and cameras omits other hardware, assembly and validation costs.
Neither project should be assigned a universal price advantage or a proven future club-face capability from a parts list or a patent analogy.

\begin{implication}
Preserve device, firmware, enabled features, ball and club preparation, installation, reference points, frames and event times.
Assess changes against repeatability, reference uncertainty and model discrepancy.
A changed face angle, spin or carry may otherwise reflect the measurement or prediction process rather than the golfer's movement.
\end{implication}
